\documentclass[11pt]{article}

% =========================
% AutoFLUKA Technical Report Template (LaTeX)
% =========================

\usepackage[margin=1in]{geometry}
\usepackage[T1]{fontenc}
\usepackage[utf8]{inputenc}
\usepackage{lmodern}
\usepackage{microtype}
\usepackage{setspace}
\usepackage{graphicx}
\usepackage{float}
\usepackage{booktabs}
\usepackage{longtable}
\usepackage{array}
\usepackage{tabularx}
\usepackage{caption}
\usepackage{subcaption}
\usepackage{amsmath,amssymb}
\usepackage{siunitx}
\usepackage{xcolor}
\usepackage{hyperref}
\usepackage{enumitem}
\usepackage{fancyhdr}
\usepackage{lastpage}
\usepackage{titlesec}
\usepackage{placeins}
\usepackage{listings}
\usepackage{etoolbox}

\hypersetup{
    colorlinks=true,
    linkcolor=blue,
    citecolor=blue,
    urlcolor=blue,
    pdftitle={AutoFLUKA Technical Report},
    pdfauthor={AutoFLUKA}
}

\setstretch{1.08}
\setlength{\parskip}{0.5em}
\setlength{\parindent}{0pt}
\setlength{\headheight}{14pt}
\captionsetup{font=small,labelfont=bf}
\sisetup{round-mode=places,round-precision=3}

\titleformat{\section}{\large\bfseries}{\thesection.}{0.5em}{}
\titleformat{\subsection}{\normalsize\bfseries}{\thesubsection.}{0.5em}{}
\titleformat{\subsubsection}{\normalsize\itshape}{\thesubsubsection.}{0.5em}{}

\pagestyle{fancy}
\fancyhf{}
\lhead{AutoFLUKA Technical Report}
\rhead{\ReportShortTitle}
\cfoot{Page \thepage\ of \pageref{LastPage}}

% =========================
% User-editable metadata
% =========================
\newcommand{\ReportTitle}{AutoFLUKA Technical Report Title}
\newcommand{\ReportShortTitle}{Case Summary}
\newcommand{\ProjectName}{Project / Campaign Name}
\newcommand{\CaseName}{Case Identifier}
\newcommand{\AuthorName}{Author / Analyst}
\newcommand{\InstitutionName}{Institution / Laboratory}
\newcommand{\ReportDate}{\today}
\newcommand{\ReportVersion}{v1.0}
\newcommand{\SoftwareName}{FLUKA / AutoFLUKA}
\newcommand{\SoftwareVersion}{Specify version if known}
\newcommand{\WorkingFolderTag}{run-folder-name}

% =========================
% Helper commands
% =========================
\newcommand{\placeholder}[1]{\textcolor{red}{[#1]}}
\newcommand{\code}[1]{\texttt{#1}}
\newcommand{\figref}[1]{Figure~\ref{#1}}
\newcommand{\tabref}[1]{Table~\ref{#1}}
\newcommand{\secref}[1]{Section~\ref{#1}}
\newcommand{\unitfmt}[1]{\,\mathrm{#1}}

\newcolumntype{Y}{>{\raggedright\arraybackslash}X}

% Allow compilation even when placeholder image files are not yet provided
\setkeys{Gin}{draft}
\makeatletter
\let\AutoFLUKAOriginalIncludegraphics\includegraphics
\renewcommand{\includegraphics}[2][]{%
  \IfFileExists{#2}{%
    \AutoFLUKAOriginalIncludegraphics[#1]{#2}%
  }{%
    \fbox{\begin{minipage}[c][0.18\textheight][c]{0.92\linewidth}\centering\small Missing image: \detokenize{#2}\end{minipage}}%
  }%
}
\makeatother

\lstset{
  basicstyle=\ttfamily\small,
  breaklines=true,
  frame=single,
  columns=fullflexible,
  showstringspaces=false,
  keywordstyle=\color{blue},
  commentstyle=\color{teal!70!black},
  rulecolor=\color{black!30}
}

% =========================
% Document start (Title page + Abstract are NOT numbered sections)
% =========================
\begin{document}

\begin{titlepage}
    \hypersetup{pageanchor=false}
    \centering
    {\Large \InstitutionName \par}
    \vspace{1.5cm}
    {\LARGE\bfseries \ReportTitle \par}
    \vspace{0.6cm}
    {\large AutoFLUKA technical report (auto-populated template) \par}
    \vspace{1.5cm}

    \begin{tabular}{rl}
        \textbf{Project:} & \ProjectName \\
        \textbf{Case:} & \CaseName \\
        \textbf{Analyst:} & \AuthorName \\
        \textbf{Software:} & \SoftwareName \\
        \textbf{Software Version:} & \SoftwareVersion \\
        \textbf{Report Version:} & \ReportVersion \\
        \textbf{Date:} & \ReportDate \\
        \textbf{Folder Tag:} & \WorkingFolderTag \\
    \end{tabular}

    \vfill
    \begin{minipage}{0.9\textwidth}
    \small
    \textbf{Template intent.} This template is designed for Monte Carlo radiation-transport campaigns executed with FLUKA and post-processed by AutoFLUKA. Replace placeholders only when required; otherwise allow AutoFLUKA to populate sections deterministically from run artifacts.
    \end{minipage}
\end{titlepage}
\hypersetup{pageanchor=true}

\begin{abstract}
\placeholder{Abstract: brief summary of the problem setup, methodology, main findings, and significance.}
\end{abstract}

\tableofcontents
\newpage

% =========================
\section{Introduction and Background}
\label{sec:intro}
\placeholder{Background context for the application area and why FLUKA/Monte Carlo is used.}

\subsection{Problem Statement and Objective}
\label{sec:problem_objective}
\begin{itemize}[leftmargin=1.5em]
  \item \textbf{Application area:} \placeholder{dosimetry | shielding | activation | detector response | accelerator | medical physics | other}
  \item \textbf{Problem being studied:} \placeholder{physical or engineering problem}
  \item \textbf{Objective:} \placeholder{what question the simulation answers}
  \item \textbf{Scope and limitations:} \placeholder{main assumptions and boundaries}
\end{itemize}

% =========================
\section{Project / Campaign Overview}
\label{sec:campaign}
\placeholder{Describe the campaign structure (smoke test vs production runs, copies/seeds), what changed across runs, and where outputs live (autofluka\_results/, Decrypted/).}

% =========================
\section{Simulation Methodology}
\label{sec:method}

\subsection{Code and execution context}
\begin{itemize}[leftmargin=1.5em]
  \item \textbf{Code:} FLUKA
  \item \textbf{FLUKA version:} \placeholder{if known}
  \item \textbf{Execution environment:} \placeholder{local | docker | wsl | other}
  \item \textbf{Workflow summary:} \placeholder{AutoFLUKA workflow summary}
\end{itemize}

\subsection{Model definition (geometry, materials, source, physics)}
\placeholder{Concise geometry/materials summary, source/beam conditions, and key physics/transport settings.}

\subsection{Scoring plan and statistical strategy}
\begin{table}[H]
\centering
\caption{Scoring plan and quantities of interest.}
\label{tab:scoring_plan}
\begin{tabularx}{\textwidth}{l Y c Y}
\toprule
\textbf{Scoring card / output} & \textbf{Quantity} & \textbf{Units} & \textbf{Reason for inclusion} \\
\midrule
\placeholder{USRBIN} & \placeholder{energy deposition map} & \placeholder{GeV/cm$^3$ per primary} & \placeholder{where dose/energy is deposited} \\
\placeholder{USRTRACK} & \placeholder{fluence spectrum} & \placeholder{1/cm$^2$ per primary} & \placeholder{spectrum at region of interest} \\
\bottomrule
\end{tabularx}
\end{table}

\begin{itemize}[leftmargin=1.5em]
  \item \textbf{Number of primaries:} \placeholder{START / NPS}
  \item \textbf{Number of runs/copies:} \placeholder{N}
  \item \textbf{Uncertainty target (if used):} \placeholder{target}
\end{itemize}

\subsection{FLUKA input summary}
\begin{itemize}[leftmargin=1.5em]
  \item \textbf{Canonical input file:} \code{\placeholder{input\_filename}}
  \item \textbf{Key regions/material assignments:} \placeholder{summary}
  \item \textbf{Important non-default cards:} \placeholder{summary}
  \item \textbf{Reproducibility note:} \placeholder{what must remain unchanged}
\end{itemize}

\subsection{Post-processing workflow}
\begin{itemize}[leftmargin=1.5em]
  \item \textbf{Decryption:} \placeholder{summary; mention Decrypted/}
  \item \textbf{Structured data products:} \placeholder{fluka\_data.json, detector\_index.json, sidecars}
  \item \textbf{Derived analyses:} \placeholder{plots, slices, isotope summaries}
\end{itemize}

% =========================
\section{Troubleshooting and Repairs}
\label{sec:troubleshooting}
\placeholder{Summarize failures, signatures, root causes, and applied fixes. Prefer the root autofluka_run_fixes.md audit file when available.}

% =========================
\section{Results and Discussion}
\label{sec:results}

\subsection{Main numerical results}
\begin{table}[H]
\centering
\caption{Main numerical results.}
\label{tab:main_results}
\begin{tabularx}{\textwidth}{l c c Y}
\toprule
\textbf{Metric} & \textbf{Value} & \textbf{Uncertainty} & \textbf{Notes} \\
\midrule
\placeholder{metric} & \placeholder{value} & \placeholder{uncertainty} & \placeholder{note} \\
\bottomrule
\end{tabularx}
\end{table}

\subsection{Figures}
\placeholder{Insert only figures that actually exist; reference them by relative path.}

\subsection{Statistical quality}
\begin{itemize}[leftmargin=1.5em]
  \item \textbf{Average uncertainties:} \placeholder{summary}
  \item \textbf{Low-statistics detectors:} \placeholder{summary}
  \item \textbf{Reliability note:} \placeholder{summary}
\end{itemize}

% =========================
\section{Reproducibility Notes}
\label{sec:repro}
\begin{itemize}[leftmargin=1.5em]
  \item \textbf{Files required to reproduce:} \placeholder{list}
  \item \textbf{Generated artifacts used in this report:} \placeholder{list}
  \item \textbf{Important run settings:} \placeholder{settings that must match}
  \item \textbf{TODO items remaining:} \placeholder{list or none}
\end{itemize}

\subsection{Lessons learned}
\label{sec:lessons}
\begin{itemize}[leftmargin=1.5em]
  \item \placeholder{Lesson 1}
  \item \placeholder{Lesson 2}
\end{itemize}

% =========================
\section{Conclusion}
\label{sec:conclusion}
\begin{itemize}[leftmargin=1.5em]
  \item \textbf{Main takeaway:} \placeholder{main conclusion}
  \item \textbf{Was the objective achieved?} \placeholder{yes/no with reason}
  \item \textbf{Recommended next steps:} \placeholder{follow-up work}
\end{itemize}

\section*{References}
\addcontentsline{toc}{section}{References}
\begin{enumerate}[leftmargin=1.8em]
  \item \placeholder{FLUKA manual or primary reference}
  \item \placeholder{paper/report/benchmark reference if used}
\end{enumerate}

\end{document}
